% !TEX root = chapter-standalone.tex



\section{Exponents interacting with products}
\label{sec:arithmetic-exponents-products}


\subsection{Functions out of products}

Next we look at the equations $z^{x \cdot y} = (z^y)^x$ and $z^{x \cdot y} = (z^x)^y$, and the respective corresponding relationships
\begin{equation}\label{eq:iso-functions-out-of-prods}
    \setC^{\setA \cartprod \setB}  \isomorphic \pars{\setC^\setB}^\setA
\end{equation}
and
\begin{equation}
    \setC^{\setA \cartprod \setB}  \isomorphic \pars{\setC^\setA}^\setB.
\end{equation}
We'll only treat the first relationship \cref{eq:iso-functions-out-of-prods} in detail, since the latter is analogous, except that certain roles are swapped.

What might be a canonical isomorphism proving \cref{eq:iso-functions-out-of-prods}? The elements of the set $ \setC^{\setA \cartprod \setB}$ are functions of two variables, of the type
\begin{equation}\label{eq:currying-function-of-two-variables}
    \mora \colon \setA \cartprod \setB \to \setC,
\end{equation}
while elements of the set $\pars{\setC^\setB}^\setA$ are functions of the type
\begin{equation} \label{eq:currying-function-to-functions}
    \morb \colon \setA \to \setC^\setB;
\end{equation}
in other words, functions of a single variable but which evaluate, for each input, to a function of the type $\setB \to \setC$.
A canonical way to turn a function of the type \cref{eq:currying-function-of-two-variables} into a function of the type \cref{eq:currying-function-to-functions} is by ``partial evaluation'': if we think of $\mora$ as having two input slots, we can partially evaluate $\mora$ by inserting an element of $\setA$ into the first slot of $\mora$, while leaving the second slot ``open'' or ``variable''.
In other words, for any element $\ela \setin \setA$, we can create, from $\mora$, the function
\begin{equation}
    \defmapperiod{
        \mora(\ela, -)
    }{
        \setB
    }{
        \to
    }{
        \setC
    }{
        \elb
    }{
        \mora (\ela, \elb)
    }
\end{equation}
But since we can do this for any $\ela \in \setA$, we have also just described a function
\begin{equation}
    \begin{aligned}
         & \setA & \to     & \ \ \setC^\setB, \\
         & \ela  & \mapsto & \ \  \mora ( \ela, -),
    \end{aligned}
\end{equation}
or, in other words, an element of $\pars{\setC^\setB}^\setA$. Since this recipe works for any function $\mora \colon \setA \cartprod \setB \to \setC$, we also now have the following function
\begin{equation}\label{eq:left-curry-function}
    \defmapcomma{
        \curry
    }{
        \setC^{\setA \cartprod \setB}
    }{
        \to
    }{
        \pars{\setC^\setB}^\setA
    }{
        \mora
    }{
        (\ela \mapsto \mora ( \ela, -))
    }
\end{equation}
which we call $\curry$.
This function has the desired type to be an isomorphism proving that $\setC^{\setA \cartprod \setB} \isomorphic \pars{\setC^\setB}^\setA$.

\begin{exercise}\label{ex:currying-is-an-iso}
    Prove that $\curry \colon \setC^{\setA \cartprod \setB} \to \pars{\setC^\setB}^\setA$ is an isomorphism.
\end{exercise}

\begin{solution}
    \todotext{JL: write up}
\end{solution}

The name ``$\curry$'' stands for ``Curry'', which is the last name of the mathematician and logician Haskell Curry (\cref{fig:haskell-curry}) who did not discover this operation, but likely made it well-known.
It is after his first name that the functional programming language Haskell is named.

The operation given by the function \cref{eq:left-curry-function} is often called ``currying''.
When needed for clarity, we call this operation \emph{right currying}, in order to distinguish it from the analogous operation of \emph{left currying} given by the function
\begin{equation}\label{eq:right-curry-function}
    \defmapperiod{
        \curry
    }{
        \setC^{\setA \cartprod \setB}
    }{
        \to
    }{
        \pars{\setC^\setA}^\setB
    }{
        \mora
    }{
        (\elb \mapsto \mora ( - , \elb))
    }
\end{equation}
We denote the functions \cref{eq:left-curry-function} and \cref{eq:right-curry-function} both by the same name ``$\curry$'' because we trust that it will be clear from context which of the two functions is being used in any given situation.

\begin{marginfigure}
    \includegraphics[width=3cm]{haskell}
    \caption{Haskell Curry (1900-1982)}
    \label{fig:haskell-curry}
\end{marginfigure}

\todotext{There is a remark (which is commented out in latex) here that could be partially included and useful perhaps, if edited.}

%\begin{remark}
%In our discussion of currying above, we started with a function $\mora \colon \setB \cartprod \setA \to \setC$ and partially evaluated it in the first variable, and left the second variable still ``open''. Alternatively, though, we could have partially evaluated $\mora$ in the \emph{second} variable, and left the first variable open, thus producing a function $\mora(-, \elb) \colon \setA \to \setC$ and leading, in a way analogous to the above, to an alternative currying operation 
%\begin{equation}
%\setC^{\setA \cartprod \setB} \to \pars{\setC^\setA}^\setB.
%\end{equation}
%In order to distinguish this operation from the currying operation previously defined, we might call this ``co-currying'', or make a distinction between ``left'' and ``right'' currying. Often, mathematicians and computer scientists follow the convention that currying is always done by partially evaluating in the first variable. This is perhaps justified by the fact that the ``co-currying'' isomorphism
%\begin{equation}
%\setC^{\setA \cartprod \setB} \to \pars{\setC^\setA}^\setB
%\end{equation}
%arising from partial evaluation in the second variable can be expressed as a composite
%\begin{equation}
%\setC^{\setA \cartprod \setB} \overset{{\setC^{\braiding}}}{\longrightarrow} \setC^{\setB \cartprod \setA} \overset{\curry}{\longrightarrow} \pars{\setC^\setA}^\setB
%\end{equation}
%of the exponentiation of $\braiding \colon \setB \cartprod \setA \to \setA \cartprod \setB$ composed with $\curry \colon \setC^{\setB \cartprod \setA} \to \pars{\setC^\setA}^\setB$.  
%\end{remark}

\subsection{Functions into products}

Now let's consider the analogue of the equation
\begin{equation}
    (x \cdot y)^z = x^z \cdot y^z
\end{equation}
for natural numbers, namely
\begin{equation}
    (\setA \cartprod \setB)^{\setC}  \isomorphic  \setA^\setC \cartprod \setB^\setC.
\end{equation}
In order to describe a canonical isomorphism which is a proof of this statement, we make an important observation about the cartesian product of sets.
Namely, for any sets $\setA$ and $\setB$, we always have two ``projection functions'' from the cartesian product $\setA \cartprod \setB$ to $\setA$ and $\setB$, respectively.
These are defined, respectively, by
\begin{equation}
    \begin{aligned}
        \proj_1 \colon & \setA \cartprod \setB & \to     & \ \ \setA, \\
                       & \tup{\ela, \elb}      & \mapsto & \ \  \ela,
    \end{aligned}
\end{equation}
and
\begin{equation}
    \begin{aligned}
        \proj_2 \colon & \setA \cartprod \setB & \to     & \ \ \setB, \\
                       & \tup{\ela, \elb}      & \mapsto & \ \  \elb,
    \end{aligned}
\end{equation}
and are called the \emph{canonical projections} associated to the cartesian product.

With the help of these projections, we define a canonical isomorphism
\begin{equation}
    \begin{aligned}
        \mora \colon & (\setA \cartprod \setB)^{\setC} & \to     & \ \  \setA^\setC \cartprod \setB^\setC, \\
                     & \phi                            & \mapsto & \ \  \tup{\phi \mthen \proj_1, \phi \mthen \proj_2}.
    \end{aligned}
\end{equation}
Illustrated diagrammatically, what $\mora$ does is, given a function
\begin{equation}
    \begin{tikzcd}
         & \setC \arrow[d, "\phi"'] & \\
         & \setA \cartprod \setB    &
    \end{tikzcd}
\end{equation}
it maps it to the pair of functions $\phi \mthen \proj_1$ and $\phi \mthen \proj_2$
\begin{equation}
    \begin{tikzcd}
              & \setC \arrow[dl, "\phi \mthen \proj_1"']   \arrow[d, "\phi"'] \arrow[dr, "\phi \mthen \proj_2"] & \\
        \setA & \setA \cartprod \setB  \arrow[l, "\proj_1"] \arrow[r, "\proj_2"']                               & \setB
    \end{tikzcd}
\end{equation}
obtained by post-composing $\phi$ with the projections $\proj_1$ and $\proj_2$, respectively.

The inverse to $\mora$ is the function $\morb \colon  \setA^\setC \cartprod \setB^\setC  \to (\setA \cartprod \setB)^{\setC}$ which takes a pair of functions $\tup{\psi_1, \psi _2} \setin \setA^\setC \cartprod \setB^\setC$ and maps it to the element of $(\setA \cartprod \setB)^{\setC}$ which is the function
\begin{equation}
    \begin{aligned}
         & \setC & \to     & \ \ \setA \cartprod \setB, \\
         & \elc  & \mapsto & \ \ \tup{\psi_1(\elc), \psi_2(\elc)}.
    \end{aligned}
\end{equation}

\begin{gradedexercise}[\exname{CanIsoFunctionsIntoProducts}]
    Check that $\mora$ and $\morb$ are indeed mutually inverse.
\end{gradedexercise}

\solutionof{CanIsoFunctionsIntoProducts}



